\section{Conclusion}

\subsection{Project Summary}
In the context of modern software development, maintaining technical documentation that is always synchronized with source code is a major challenge for development teams. Many projects face the issue of documentation quickly becoming outdated due to infrequent updates, causing difficulties in maintenance, expansion, and system handover.

The \textbf{LivingDocs} project was proposed to address this problem by combining modern technologies such as GitHub Integration, Artificial Intelligence (AI), Retrieval-Augmented Generation (RAG), and Workflow Automation. The system aims to automate the creation, updating, and management of technical documentation based on source code changes, thereby significantly reducing manual workload and improving documentation quality.

Within the scope of the Software Engineering course, the team conducted requirement analysis, designed the overall architecture, and developed descriptive documentation for the main system components. These deliverables form the foundation for future development and deployment phases.

\subsection{Project Achievements}
During the project, the team accomplished the following key tasks:
\begin{itemize}
    \item Analyzed functional and non-functional requirements of the system.
    \item Designed GitHub integration workflow using OAuth2 and GitHub API.
    \item Built an automated documentation update process triggered by GitHub events.
    \item Proposed an AI architecture to support source code analysis and technical documentation generation.
    \item Designed a version management mechanism for documentation and detection of inconsistencies between documentation and source code.
    \item Developed testing procedures and a system quality evaluation plan.
\end{itemize}

These results contribute to forming a modern documentation management model that effectively supports software development in a systematic way.

\subsection{Project Limitations}
Due to time constraints and the scope of the course, the project still has certain limitations:
\begin{itemize}
    \item The system remains at the design and solution description stage, without full implementation of all functions.
    \item Some components such as the AI Engine, Vector Database, and Workflow Automation have not been tested in real environments.
    \item Performance, scalability, and security evaluations have not yet been conducted.
    \item Full integration with CI/CD processes in real development environments has not been implemented.
\end{itemize}

These limitations will serve as the basis for further improvement in subsequent development phases.

\subsection{Future Development}
In the future, the LivingDocs system can be expanded in various directions to enhance usability and practical application. Potential development directions include:
\begin{itemize}
    \item Full implementation of Backend, Frontend, and Database components.
    \item Integration of AI models with more accurate source code analysis capabilities.
    \item Extending support to multiple source code management platforms such as GitLab or Bitbucket.
    \item Adding automated testing and continuous deployment (CI/CD) mechanisms.
    \item Developing a more detailed role-based access control system.
    \item Optimizing system performance for large-scale projects.
    \item Strengthening security and monitoring solutions during operation.
\end{itemize}

Advancing along these directions will help LivingDocs become an intelligent technical documentation management platform, better meeting the needs of real-world software development teams.

\subsection{Conclusion}
Overall, the \textbf{LivingDocs} project has proposed a solution to support technical documentation management in an automated, intelligent, and synchronized manner with the software development process. Through GitHub integration, AI application, and workflow automation, the system aims to reduce manual work, improve documentation accuracy, and enhance collaboration efficiency within development teams.

Although the project is currently at the design and documentation stage, the achieved results demonstrate the feasibility of the solution and lay the groundwork for future implementation. With the proposed development directions, LivingDocs has the potential to become a modern documentation management tool, contributing to improved quality in software development and maintenance in real-world projects.
